% ===========================================================================
%  Character-Level Language Model for Robot Instruction Following
%  Nature-style manuscript (compatible with nature.cls or article class)
% ===========================================================================

%\documentclass{nature}
\documentclass[11pt,a4paper]{article}

% --- Packages ---
\usepackage[utf8]{inputenc}
\usepackage[T1]{fontenc}
\usepackage{newtxtext,newtxmath}
\usepackage{graphicx}
\usepackage{booktabs}
\usepackage{amsmath,amssymb}
\usepackage{hyperref}
\usepackage[margin=2.5cm]{geometry}
\usepackage{subcaption}
\usepackage{xcolor}
\usepackage{listings}
\usepackage{caption}
\usepackage{enumitem}
\usepackage{float}

% --- Hyperref setup ---
\hypersetup{
    colorlinks=true,
    linkcolor=blue,
    urlcolor=blue
}

% --- Code listing style ---
\lstset{
    basicstyle=\ttfamily\small,
    breaklines=true,
    frame=single,
    numbers=left,
    numberstyle=\tiny,
    xleftmargin=2em,
    language=Python,
    showstringspaces=false,
    columns=flexible
}

% --- Title ---
\title{A Character-Level Transformer for Mapping Natural Language Instructions to Structured Robot Commands}

\author{Author Name}

\date{}

\begin{document}

\maketitle

% ===========================================================================
%  ABSTRACT
% ===========================================================================
\begin{abstract}
Translating natural language instructions into structured robot control commands requires models that are both syntactically precise and data-efficient.
Existing large language models achieve strong performance on general-purpose instruction following, yet remain over-parameterized for domain-specific structured-output tasks,
where the output space is constrained to a small set of JSON templates.
Here we train a 3.2-million-parameter GPT-2-style Transformer with character-level tokenization on 887 robot instruction--command pairs
and evaluate it on 50 held-out test instructions under deterministic greedy decoding.
The model produces syntactically valid JSON in 100\% of test cases, achieves 80.0\% exact-match accuracy,
and correctly predicts the target function name in all 50 samples.
Error analysis reveals that residual failures are concentrated in numerical precision and cross-unit conversion rather than structural errors.
These results suggest that character-level encoding, when paired with careful data normalization,
offers a compact and reliable route for structured-output generation in resource-constrained settings.
\end{abstract}

% ===========================================================================
%  INTRODUCTION
% ===========================================================================
\section{Introduction}

Translating natural language (NL) instructions into executable robot commands is a core capability for human--robot interaction systems.
The task requires mapping free-form text---``Move the TCP to (0.5, 0.3, 0.7)~m''---to structured JSON objects
such as \texttt{\{"function": "move\_tcp", "kwargs": \{"x": 500.0, "y": 300.0, "z": 700.0\}\}}.
Two conflicting demands define the problem: the output must be \emph{syntactically exact} (a single misplaced bracket breaks downstream execution),
yet the input is \emph{linguistically variable} (users express identical intent with diverse phrasing and units).

Large pre-trained language models perform well on instruction-following benchmarks but incur substantial inference cost
and require tokenizers whose subword splits can introduce unpredictable artifacts in JSON output.
For domain-constrained settings---where the command vocabulary is small (fewer than 10 function types)
and the output syntax is fixed---a character-level autoregressive model may be sufficient.
The nanoGPT framework provides a minimal, auditable GPT-2 implementation
that can be rapidly adapted to such tasks.

In this work, we train and evaluate a small Transformer ($\sim$3.2M parameters) on the Studeni/robot-instructions dataset,
using character-level encoding with a vocabulary of only 77 tokens.
We report four findings:
(1)~character-level tokenization eliminates JSON syntax errors entirely (100\% validity on test);
(2)~80.0\% of generated commands match the reference exactly after canonicalization;
(3)~function-name selection is perfect across all test cases; and
(4)~remaining errors are concentrated in numerical precision and unit conversion, not in structural reasoning.

% ===========================================================================
%  RESULTS
% ===========================================================================
\section{Results}

\subsection{Training Dynamics}

The model was trained for 3,000 iterations with a batch size of 32 and block size of 256,
processing approximately 24.6 million tokens in total.
Training and validation loss decreased smoothly under a cosine learning-rate schedule with 50-step linear warmup (Fig.~\ref{fig:training}).
The final training loss reached 0.087 and validation loss 0.094, corresponding to perplexities of 1.09 and 1.10, respectively.
Model FLOPs utilization (MFU) stabilized at approximately 28\% on a single consumer GPU,
consistent with the small model size and low arithmetic intensity of character-level sequences.

\begin{figure}[H]
    \centering
    \includegraphics[width=\columnwidth]{out-robot-instr/training_metrics.png}
    \caption{Training metrics over 3,000 iterations.
    \textbf{(a)} Train and validation loss.
    \textbf{(b)} Raw per-step loss.
    \textbf{(c)} Perplexity ($e^{\text{loss}}$).
    \textbf{(d)} Learning-rate schedule (warmup + cosine decay).
    \textbf{(e)} Model FLOPs utilization (MFU).
    \textbf{(f)} Summary statistics.}
    \label{fig:training}
\end{figure}

\subsection{Evaluation on Held-Out Instructions}

We evaluated the final checkpoint on 50 held-out test instructions using greedy decoding (temperature = 0)
with a maximum of 500 new tokens and a stop condition at the double-newline delimiter.
Four metrics were computed (Table~\ref{tab:metrics}):

\begin{table}[H]
    \centering
    \caption{Aggregate evaluation metrics on 50 test instructions.}
    \label{tab:metrics}
    \begin{tabular}{@{}lcc@{}}
        \toprule
        \textbf{Metric} & \textbf{Value} & \textbf{Samples} \\
        \midrule
        JSON validity        & 100.0\% & 50/50 \\
        Exact match          &  80.0\% & 40/50 \\
        Function-name match  & 100.0\% & 50/50 \\
        Character accuracy   & $\approx$\,93.8\% & --- \\
        \bottomrule
    \end{tabular}
\end{table}

The JSON validity rate of 100\% indicates that character-level encoding provides a strong inductive bias toward syntactically correct output:
every bracket, comma, and quotation mark is modeled as an independent character,
leaving no opportunity for subword-merging artifacts that can corrupt structured output in BPE-based tokenizers.

\subsection{Error Analysis}

The 10 samples that did not achieve exact match were manually inspected and classified into three error types:

\paragraph{Type I: Numerical deviation (most common).}
The model generated the correct function name and JSON structure but produced incorrect numerical values.
Examples include joint-index substitution (\texttt{joint [1]} $\to$ \texttt{[3]}) and coordinate drift (\texttt{z=300.0} $\to$ \texttt{z=350.0}).

\paragraph{Type II: Unit conversion error.}
When the instruction specified coordinates in centimeters (``Move to (100, --50, 200)~cm''),
the expected output required millimeter conversion ($\times 100$). The model produced values closer to the raw input,
indicating that unit-conversion reasoning is not reliably acquired from the training distribution alone.

\paragraph{Type III: Multi-command merge/split error.}
Instructions requiring multiple independent robot commands were sometimes incorrectly merged into a single command,
or a single multi-joint command was split across separate calls.

These failure modes suggest that the primary bottleneck is not structural understanding but \emph{numerical reasoning} under sparse training coverage.
Augmenting the training data with explicit unit-conversion examples or adding a symbolic preprocessing step
for unit normalization could address the majority of remaining errors.

% ===========================================================================
%  DISCUSSION
% ===========================================================================
\section{Discussion}

The results demonstrate that a character-level Transformer with fewer than 4 million parameters
can serve as a reliable structured-output generator for domain-specific robot instruction following.
The near-perfect JSON validity and function-name accuracy are consistent with the hypothesis
that character-level modeling removes the tokenization-related failure modes that affect subword-based models on structured tasks.

The 80\% exact-match rate must be interpreted in context: the remaining 20\% of errors are dominated by numerical imprecision,
not by incorrect command structure or function selection.
This pattern suggests that the model has successfully learned the \emph{syntax} and \emph{semantics} of the command space,
and that residual errors reflect insufficient coverage of arithmetic transformations rather than a fundamental modeling limitation.

Several limitations should be noted.
First, the dataset is small (887 training samples) and covers only three function types;
generalization to larger command vocabularies or more complex JSON schemas remains untested.
Second, the model uses absolute positional encodings with a fixed context window of 256 tokens,
which may limit performance on multi-step instruction sequences.
Third, the evaluation was conducted with greedy decoding (temperature = 0);
stochastic decoding strategies may yield different trade-offs between diversity and precision.

Future work could explore data-augmentation strategies that explicitly target unit-conversion reasoning,
integration of relative positional encodings (e.g., RoPE) for longer sequences,
and comparison with instruction-tuned models of comparable size.

% ===========================================================================
%  METHODS
% ===========================================================================
\section{Methods}

\subsection{Data Preprocessing}

The dataset Studeni/robot-instructions contains natural-language-to-JSON command pairs
for robotic manipulation tasks.
We developed a preprocessing pipeline (\texttt{data/robot\_instr/prepare.py}) that performs the following steps:

\begin{enumerate}[leftmargin=*,nosep]
    \item \textbf{Loading}: data are read from a local \texttt{input.txt} cache (986 formatted samples),
          or downloaded from HuggingFace Hub on first use.
    \item \textbf{Cleaning}: all floating-point numbers are rounded to two decimal places via \texttt{round\_floats()};
          JSON outputs are canonicalized through \texttt{json.loads} $\to$ \texttt{json.dumps}.
    \item \textbf{Formatting}: each sample is serialized as
          \texttt{Human: \textless{}instruction\textgreater{}\textbackslash nRobot: \textless{}json\textgreater\textbackslash n\textbackslash n}.
    \item \textbf{Splitting}: a sequential 9:1 train/validation split yields 887 training and 99 validation samples.
    \item \textbf{Encoding}: a character-level vocabulary is built from the union of all characters in the training and validation sets.
\end{enumerate}

The preprocessing addresses two data-quality issues in the raw dataset:
(i)~inconsistent distance units (meters, centimeters, millimeters) across instructions, and
(ii)~unconstrained decimal precision in angular values (e.g., $\pi/4$ vs.\ integer degrees).

\subsection{Character-Level Tokenization}

We adopt a character-level tokenizer following the \texttt{shakespeare\_char} recipe in nanoGPT.
The vocabulary consists of $V = 77$ unique characters, comprising uppercase and lowercase letters, digits,
JSON syntax tokens (\texttt{\{\}[]",:}), the minus sign, decimal point, whitespace, and newline characters.
Bidirectional mapping is stored in \texttt{meta.pkl}:

\begin{equation}
    \text{encode}(s) = [\,\text{stoi}[c] \mid c \in s\,], \qquad
    \text{decode}(\mathbf{ids}) = \text{concat}([\,\text{itos}[i] \mid i \in \mathbf{ids}\,])
\end{equation}

Token IDs are serialized as \texttt{uint16} arrays in \texttt{train.bin} (240~KB) and \texttt{val.bin} (26~KB).

The choice of character-level encoding is motivated by three considerations:
(1)~the tiny vocabulary (77 tokens) minimizes embedding parameters ($77 \times 256 = 19{,}712$ weights);
(2)~JSON syntax precision is guaranteed, since every bracket and comma is a distinct token; and
(3)~out-of-vocabulary tokens cannot occur by construction.

\subsection{Model Architecture}

The model follows the GPT-2 architecture implemented in nanoGPT,
with hyperparameters configured for the small-data regime (Table~\ref{tab:arch}):

\begin{table}[H]
    \centering
    \caption{Model architecture configuration.}
    \label{tab:arch}
    \begin{tabular}{@{}lc@{}}
        \toprule
        \textbf{Hyperparameter} & \textbf{Value} \\
        \midrule
        Transformer layers ($L$)     & 4 \\
        Attention heads ($H$)        & 4 \\
        Embedding dimension ($d$)    & 256 \\
        Maximum sequence length ($T$) & 256 \\
        Dropout rate                 & 0.2 \\
        Bias parameters              & None \\
        Total parameters             & $\sim$3.2M \\
        \bottomrule
    \end{tabular}
\end{table}

The forward pass is defined as:

\begin{align}
    \mathbf{h}_0 &= \text{Dropout}\big(\text{wte}(x) + \text{wpe}(\text{pos})\big) \\
    \mathbf{h}_\ell &= \mathbf{h}_{\ell-1} + \text{MLP}_\ell\big(\text{LN}_{\ell,2}\big(
        \mathbf{h}_{\ell-1} + \text{Attn}_\ell(\text{LN}_{\ell,1}(\mathbf{h}_{\ell-1}))
    \big)\big), \quad \ell = 1,\dots,L \\
    P(y \mid x) &= \text{softmax}\big(\text{LN}_f(\mathbf{h}_L) \cdot \mathbf{W}_e^\top\big)
\end{align}

where $\text{wte}$ and $\text{lm\_head}$ share weights ($\mathbf{W}_e^\top$),
causal self-attention uses Flash Attention when PyTorch $\geq$ 2.0,
the MLP block uses a 4$\times$ expansion ratio with GELU activation,
and residual projections are initialized with a $1/\sqrt{2L}$ scaling factor for training stability.

\subsection{Training Configuration}

Training uses the AdamW optimizer with $\beta_1 = 0.9$, $\beta_2 = 0.99$,
weight decay $\lambda = 0.1$ (applied only to parameters with $\geq 2$ dimensions),
and gradient clipping at $1.0$.
The learning rate follows a warmup-then-decay schedule:

\begin{equation}
    \eta(t) =
    \begin{cases}
        \eta_{\max} \cdot \frac{t}{t_{\text{warm}}}, & t \leq t_{\text{warm}} \\[4pt]
        \eta_{\min} + \frac{\eta_{\max} - \eta_{\min}}{2}
            \left(1 + \cos\left(\pi \cdot \frac{t - t_{\text{warm}}}{T - t_{\text{warm}}}\right)\right), & t > t_{\text{warm}}
    \end{cases}
\end{equation}

with $\eta_{\max} = 10^{-3}$, $\eta_{\min} = 10^{-4}$, $t_{\text{warm}} = 50$, and $T = 3000$ iterations.
Mixed-precision training (\texttt{bfloat16} autocast) and asynchronous data prefetching
(\texttt{pin\_memory} + \texttt{non\_blocking}) are used to reduce memory footprint and overlap data transfer with computation.
Checkpoints are saved only when validation loss improves, implementing a lightweight form of early stopping.

Throughput per iteration is:

\begin{equation}
    \text{tokens/iter} = g_{\text{acc}} \times B \times T = 1 \times 32 \times 256 = 8{,}192,
\end{equation}

yielding $8{,}192 \times 3{,}000 = 24{,}576{,}000$ total training tokens.

\subsection{Evaluation Protocol}

Evaluation is performed using \texttt{sample.py} in evaluation mode on 50 held-out test instructions.
Decoding uses greedy selection (temperature = 0) with a stop condition at the double-newline token or after 500 new tokens.
Generated outputs are compared against reference commands using four metrics:

\begin{itemize}[leftmargin=*,nosep]
    \item \textbf{JSON validity}: whether the output parses as valid JSON.
    \item \textbf{Exact match}: whether the canonicalized JSON (sorted keys, normalized whitespace) matches the reference.
    \item \textbf{Function-name match}: whether the set of function names in the generated output matches the reference.
    \item \textbf{Character error rate} (CER): Levenshtein distance between generated and reference strings, normalized by reference length.
\end{itemize}

% ===========================================================================
%  CONCLUSION
% ===========================================================================
\section{Conclusion}

We have shown that a 3.2M-parameter character-level Transformer,
trained on 887 robot instruction--command pairs,
generates syntactically valid JSON in 100\% of test cases and achieves 80\% exact-match accuracy.
The model correctly identifies the target function in every test sample,
and residual errors are confined to numerical precision rather than structural understanding.
These findings indicate that character-level encoding is a practical design choice
for domain-specific structured-output tasks where the output vocabulary is small and precision is paramount.
Data augmentation targeting unit conversion and multi-step composition represents a natural next step
toward closing the remaining performance gap.

% ===========================================================================
%  DATA AND CODE AVAILABILITY
% ===========================================================================
\section*{Data and Code Availability}

The nanoGPT framework is available at \url{https://github.com/karpathy/nanoGPT}.
The Studeni/robot-instructions dataset is available on HuggingFace Hub.
Preprocessing scripts, model configurations, and evaluation code are included in the project repository.

\end{document}
